% Exemplo: glossário pós-textual simples (termos técnicos com definições, em ordem alfabética)
% Compilar: latexmk -pdf glossario-manual.tex     Testado com: abntex2, TeX Live 2024 ou mais recente
% Para glossários grandes, com remissões e referências cruzadas, veja o pacote glossaries-extra.
\documentclass[12pt, a4paper, brazilian]{abntex2}
\usepackage[T1]{fontenc}

%%BEGIN:glossariodef
% Ambiente "glossario": título em maiúsculas (entra no sumário) e um
% parágrafo por termo
\newcommand{\termo}[2]{\par\noindent\textbf{#1:} #2\par\medskip}
\newenvironment{glossario}{%
  \pretextualchapter{\MakeUppercase{Glossário}}%
  \addcontentsline{toc}{chapter}{\MakeUppercase{Glossário}}%
  \begingroup\SingleSpacing}{\endgroup}
%%END:glossariodef

\begin{document}

\tableofcontents*

\textual
\chapter{Introdução}

Texto do trabalho.

\postextual

%%BEGIN:glossariouso
\begin{glossario}
  \termo{API}{Conjunto de definições e protocolos que permite a comunicação
    entre sistemas.}
  \termo{Framework}{Estrutura reutilizável de software que oferece
    funcionalidades genéricas.}
  \termo{SGBD}{Sistema que gerencia a criação, o acesso e a manutenção de
    bancos de dados.}
\end{glossario}
%%END:glossariouso

\end{document}
